%% ============================================================================
\section{Materials and Datasets}
\label{ch:materials}
\addcontentsline{lot}{section}{Materials and Datasets}
%% ============================================================================

\subsection{Overview}

This chapter describes the data and the model architectures used throughout the
thesis.
The experimental subject is a semantic segmentation workload over satellite imagery
of the southern Mesopotamian floodplain, and the object of study is the data
pipeline that feeds that workload on a shared HPC filesystem.
The dataset is therefore treated in two distinct ways: as a machine learning
dataset, characterised by its imagery, class balance and evaluation metrics
(Section~\ref{sec:dataset}); and as an I/O workload, characterised by file count,
file size distribution and access pattern (Section~\ref{sec:io-workload}).
The second characterisation is the one that motivates \pads{}, and it is the one
usually left implicit in the segmentation literature.

\subsection{Dataset: Casini et al. (2023)}
\label{sec:dataset}

The main experiments use the dataset introduced by Casini et al.~\cite{Casini23},
which addresses semantic segmentation and detection of mounded archaeological sites
- tells - across the southern Mesopotamian floodplain in Iraq.

\subsubsection{Provenance and construction}

The dataset originates from a corpus of geo-referenced vector shapes corresponding
to the contours of known mound sites within the survey area of the Floodplains
Project, spanning approximately 66{,}000~km\textsuperscript{2}.
The vector corpus was assembled at the University of Bologna by consolidating all
published archaeological surveys in the region and geo-referencing the sites
catalogued therein; it contains 4{,}934 shapes, each corresponding to a site
confirmed by ground-truthing and by study of the associated surface scatter of
artefacts.

Because the corpus was compiled as an archaeological reference rather than as a
training set, the authors applied two filtering passes before image generation:

\begin{enumerate}
  \item \textbf{Removal of the 200 largest sites by area.}
    These exceed the area of the square input region and would therefore yield an
    almost entirely positive segmentation mask, providing no usable learning signal.
    Visual inspection confirmed that they follow the shape of extended areas rather
    than of simple mounds.
  \item \textbf{Removal of 684 sites} that were either too small to be a tell or
    marked as destroyed.
    The size threshold was set at roughly 1{,}000~m\textsuperscript{2},
    corresponding to a circle of 30~m diameter.
    Sites below this threshold typically correspond to generic annotations for known
    sites of unknown size or imprecise location.
\end{enumerate}

Image generation was performed in QGIS.
For each retained site, a Python script exported a square of side $L$ centred on the
site centroid containing only Bing Maps satellite imagery, accessed through the
QuickMapService plugin.
The same square was exported a second time without the base map but with the site
contour rendered as a solid filled shape, giving the corresponding binary
ground-truth mask.
The task is therefore binary: one class for mounded sites, one for everything else.

Two spatial configurations were produced: $L = 1000$~m (``1k'') and $L = 2000$~m
(``2k'').
Imagery was exported at approximately 1 pixel per metre - 1024~px for 1000~m, and
double that for the larger extent - and subsequently downscaled by a factor of two
to reduce computational cost, which the original authors report as having low impact
on accuracy.
At training time a square of side $L/2$ is randomly cropped from the exported tile;
this both prevents the model from learning a positional prior in which sites always
appear at the image centre, and acts as data augmentation.
Further augmentation consists of random rotation, mirroring, and small shifts in
brightness and contrast, resampled at each training iteration.

Finally, 1{,}155 images with empty masks - containing no site to predict - were
added, sampled from locations proposed by the archaeologists: highly urbanised
areas, intensive agriculture, flood-prone terrain such as artificial lakes and
basins, and rocky hills and mountains.
These negatives are what force the model to learn what a tell is \emph{not}.
The final dataset consists of 5{,}025 image/mask pairs, split 90\% training /
10\% held-out test, stratified over the empty images, with a further 10\% of the
training partition randomly reserved as a validation set.

A CORONA-augmented variant also exists, in which declassified historical satellite
imagery is imported into QGIS and cropped identically to the Bing tiles, giving a
second input channel group.
Because CORONA was used only as a complement in the original work, sites destroyed
after the 1970s were excluded from that analysis.

\paragraph{Dataset variant used in this thesis.}
All experiments reported here use the Bing 1k configuration unless otherwise stated.
This choice is deliberate: the 1k configuration is the smallest of the published
variants and therefore the least favourable to a data-staging argument - a
pipeline that is I/O-bound on 1k tiles will be at least as I/O-bound on 2k tiles or
on the CORONA-augmented input.
Section~\ref{ch:ablation} revisits the larger input configurations.

\subsubsection{Class imbalance}

The segmentation target is severely imbalanced: mound pixels constitute a small
minority of any given tile, and 1{,}155 of the 5{,}025 tiles contain no positive
pixels at all.
Direct measurement over the corpus used here gives a mean positive-pixel fraction of
approximately 9.6\% per site-bearing tile, with a median near 7\%.
The masks themselves are cleanly binary - 90.4\% of mask pixels take the
background value and 8.8\% the foreground value, with only 0.7\% at intermediate
values - so label ambiguity at object boundaries is not a material source of error.

This imbalance is the reason the original work adopts focal loss~\cite{Lin17} over
Dice loss, and the reason the primary reported metric is IoU rather than pixel
accuracy.
Focal loss is one of several loss functions purpose-built for class imbalance in
segmentation; Azad et al.~\cite{Azad23} survey twenty-five such loss functions and
describe focal loss's reweighting term as reducing the loss contribution of
easily-classified samples relative to hard-to-classify ones.
In a tile that is 90\% background, as here, the great majority of easily-classified
pixels are background, so this reweighting acts as an implicit foreground
up-weighting -- though the survey states the mechanism in general easy/hard terms,
not specifically in terms of this dataset's class split.
It is also, importantly for this thesis, the reason the workload is decode-heavy
relative to its useful signal: a substantial fraction of the bytes read from the
filesystem contribute nothing but negative pixels.

\subsubsection{Evaluation metrics}

The original paper evaluates using intersection-over-union at the pixel level, and
accuracy, precision, recall and the Matthews correlation coefficient at the
site-detection level:
%
\begin{align}
  \text{IoU} &= \frac{|P \cap G|}{|P \cup G|} \\[4pt]
  \text{Accuracy} &= \frac{TP + TN}{TP + TN + FP + FN} \\[4pt]
  \text{Recall} = \frac{TP}{TP + FN}, \quad
  &\text{Precision} = \frac{TP}{TP + FP} \\[4pt]
  \text{MCC} &= \frac{TP \times TN - FP \times FN}
                    {\sqrt{(TP+FP)(TP+FN)(TN+FP)(TN+FN)}}
\end{align}
%
where $P$ denotes the predicted shape and $G$ the ground-truth shape.

The distinction between the two levels matters.
IoU measures how well a predicted contour overlaps an annotated one; it does not
measure how many sites were found.
The original authors note that IoU can improve while detection accuracy degrades, as
occurs between their Model~5 and Model~6: a more conservative model produces fewer
positive predictions, shrinking the union term and raising IoU, while missing sites
and lowering recall.

\paragraph{Reduction convention.}
A subtlety that materially affects reported numbers is how per-image scores are
aggregated.
The reference implementation reports
\code{smp.metrics.iou\_score(..., reduction="macro-imagewise")}: IoU is computed per
image and then averaged, with the library's convention that an empty ground truth
correctly predicted empty scores $1.0$.
This differs from a globally pooled IoU, in which true positives, false positives
and false negatives are summed across the whole set before the ratio is taken.
On the same trained model the two conventions differ by several percentage points,
with the per-image reduction the higher of the two because the empty tiles
contribute perfect scores.
All comparisons in this thesis therefore state which reduction is used.
Section~\ref{sec:repro-metric} quantifies the difference.

\subsubsection{Reference results}

Casini et al. report the IoU scores in Table~\ref{tab:casini-reference} on the
held-out test set.
Standard deviations arise from repeated testing to average out the random crop,
which is applied at test time as well as during training.

\begin{table}[!ht]
\centering
\begin{tabular}{llccc}
\toprule
Model & Input & Filter & IoU (\%) & Std.\ dev \\
\midrule
Model 1 & Bing 1k          &            & 74.17 & 0.38 \\
Model 2 & Bing 1k          & \checkmark & 78.10 & 0.54 \\
Model 3 & Bing + CORONA 1k &            & 74.06 & 0.39 \\
Model 4 & Bing 2k          &            & 79.77 & 0.34 \\
Model 5 & Bing 2k          & \checkmark & 81.54 & 0.35 \\
Model 6 & Bing + CORONA 2k & \checkmark & 83.45 & 0.18 \\
\bottomrule
\end{tabular}
\caption{IoU scores reported by Casini et al.~\cite{Casini23} for their
experimental configurations.}
\label{tab:casini-reference}
\end{table}

At the detection level, Model~5 achieves 62.57\% accuracy under automatic evaluation
and 80.40\% after human-in-the-loop adjustment for mislabelled (no-longer-visible)
sites; Model~6 achieves 60.08\% and 79.13\% respectively.
Their configuration search compared U-Net against MA-Net, ResNet-18 against
EfficientNet-B3, and Dice loss against focal loss, settling on
MA-Net + EfficientNet-B3 + focal loss trained for 20 epochs as the best
configuration; they observe that the differences between configurations were small,
within a few percentage points, and plausibly attributable to fluctuation from the
random augmentation.

These figures serve two purposes here.
First, they establish an accuracy floor: any optimisation to the data pipeline
proposed in this thesis must be shown not to degrade segmentation quality relative
to this reference.
Second, the observation that architectural choice moves IoU by only a few points
motivates the framing of this work - if the model is not where the remaining
headroom is, the pipeline may be.

\subsection{The dataset as an I/O workload}
\label{sec:io-workload}

The above is the standard description of the dataset.
For the purposes of this thesis a second description is required, because the
property of the dataset that \pads{} responds to is not its class balance but its
file layout.

The Bing 1k dataset consists of 5{,}025 image files and 5{,}025 corresponding mask
files stored as individual objects on the shared GPFS filesystem.
Each training step therefore issues, per sample, an independent
\code{open}/\code{stat}/\code{read}/\code{close} sequence against a parallel
filesystem that is optimised for large sequential transfers rather than for
many-small-file metadata operations.
The relevant properties are:

\begin{description}
  \item[Many small files.] Metadata cost per sample is close to constant regardless
    of how few bytes are actually read, so per-file overhead does not amortise.
  \item[Paired access.] Every image read is accompanied by a mask read at a
    different path, doubling the metadata traffic per sample.
  \item[Shared contention.] GPFS is a shared resource; observed read latency is not
    a property of the dataset alone but of concurrent jobs on the cluster, and is
    therefore non-stationary across a run.
  \item[Deterministic access order.] Once the epoch permutation is fixed, the
    sequence of samples - and hence of files - is known in advance.
    Future data access is predictable, which is precisely the precondition that
    makes proactive prefetching and shard staging exact rather than speculative.
\end{description}

The final point is the hinge of the entire method.
The dataset is small enough to fit on node-local scratch and its access pattern is
known ahead of time; the fact that the baseline pipeline nonetheless leaves the GPU
idle is a property of the loader, not of the data.

\subsection{Models}

Three segmentation architectures are used, spanning a CNN encoder--decoder, an
attention-augmented CNN, and a transformer.
All three were trained under an identical data pipeline so that pipeline effects can
be isolated from architectural effects.

\subsubsection{MA-Net}

MA-Net (Multi-scale Attention Network)~\cite{Fan20} extends the U-Net
encoder--decoder with a self-attention mechanism inserted in the latent space,
allowing the network to weight latent features according to content.
Originally developed for liver and tumour segmentation in medical imaging, it has
since been applied to remote sensing.
It is included here because it is the architecture selected as best-performing by
Casini et al., and therefore constitutes the reproduction baseline for this thesis.
Implementation: \code{segmentation\_models\_pytorch}~\cite{Iakubovskii19} with an
EfficientNet-B3~\cite{Tan19} encoder pretrained on ImageNet~\cite{Deng09}.

\subsubsection{DeepLabV3+}

DeepLabV3+~\cite{Chen18} uses atrous (dilated) convolution and atrous spatial
pyramid pooling to capture multi-scale context without loss of spatial resolution,
combined with a decoder module that recovers object boundaries.
It is included as a widely used, high-capacity CNN baseline and represents the
heaviest of the three in both parameter count and per-step GPU time.
Implementation: \code{segmentation\_models\_pytorch} with a ResNet-50 encoder.

\subsubsection{SegFormer}

SegFormer~\cite{Xie21} pairs a hierarchical, positional-encoding-free transformer
encoder with a lightweight all-MLP decoder.
It is included as the transformer baseline and, critically for this thesis, as the
least GPU-bound of the three architectures: its low per-step compute makes it the
configuration in which data-loading stalls are most exposed.
Implementation: the NVlabs reference architecture as distributed through the
HuggingFace \code{transformers} library.

\subsubsection{Measured characteristics (Bing 1k, default loader)}
\label{sec:measured-baseline}

Table~\ref{tab:baseline-default} reports measurements on the Iris cluster on a
single V100 with the baseline PyTorch \code{DataLoader} at default settings, and
Table~\ref{tab:baseline-tuned} the same with loader parameters tuned
(\code{num\_workers=4}, \code{pin\_memory=True}). These figures were recollected
in full for this thesis (Section~\ref{ch:results}, Table~\ref{tab:rq1-loader-util}
reports the same underlying runs); an earlier draft of this section used numbers
from a prior, smaller measurement campaign that this table now supersedes, since
the two disagreed once the full evaluation was collected -- notably on the
direction of SegFormer's IoU change under tuning, not only its magnitude.

\begin{table}[!ht]
\centering
\begin{tabular}{lccc}
\toprule
 & SegFormer & MA-Net & DeepLabV3+ \\
\midrule
Trainable parameters & 3.7\,M & 17.2\,M & 26.7\,M \\
Model size (MB)      & 15     & 68      & 106 \\
IoU                  & 0.3974 & 0.3440 & 0.3240 \\
Epoch reached         & 36     & 12      & 2 \\
Batch size           & 32     & 32      & 32 \\
GPU utilisation (\%) & 16.0   & 34.4    & 33.6 \\
\bottomrule
\end{tabular}
\caption{Baseline characteristics with the default PyTorch \code{DataLoader}
(\code{num\_workers=0}). IoU is the main-pipeline pooled convention
(Section~\ref{sec:repro-metric}), not the reference-comparable per-image mean.}
\label{tab:baseline-default}
\end{table}

\begin{table}[!ht]
\centering
\begin{tabular}{lccc}
\toprule
 & SegFormer & MA-Net & DeepLabV3+ \\
\midrule
Trainable parameters & 3.7\,M & 17.2\,M & 26.7\,M \\
Model size (MB)      & 15     & 68      & 106 \\
IoU                  & 0.3725 & 0.3695 & 0.3560 \\
Epoch reached         & 28     & 36      & 28 \\
Batch size           & 32     & 32      & 32 \\
GPU utilisation (\%) & 40.5   & 85.1    & 82.6 \\
\bottomrule
\end{tabular}
\caption{Characteristics with loader parameters tuned
(\code{num\_workers=4}, \code{pin\_memory=True}).}
\label{tab:baseline-tuned}
\end{table}

Two observations from these tables set up Section~\ref{ch:methodology}, and the
second is more precisely stated here than an earlier draft had it. First, GPU
utilisation moves by 24--51 percentage points between the two tables depending
on architecture -- large enough, together with the per-stage timing breakdown of
Section~\ref{sec:res-rq1}, that the data path is a plausible bottleneck rather
than a rounding error. Second, the same two loader flags do \emph{not} move IoU
by a consistent, large margin: DeepLabV3+ gains 3.2 points and MA-Net 2.6, but
SegFormer \emph{loses} 2.5 -- mixed in direction and modest in magnitude next to
the reference study's architectural search, not comparable to it. The empirical
hook for this work is therefore not that loader configuration rivals
architecture choice as a lever on accuracy; the data below does not support that
claim. It is that loader configuration is where this workload's throughput
headroom demonstrably sits, independent of what it does or does not do to IoU,
and that headroom is worth pursuing on its own terms
(Section~\ref{ch:conclusions} returns to this distinction directly).

\subsection{Hardware and software environment}

All primary experiments were run on the Iris cluster at the University of
Luxembourg, on GPU nodes equipped with NVIDIA V100 accelerators, scheduled via
Slurm.
Storage is provided by a shared GPFS filesystem, with node-local scratch available
for staging.
The software stack is containerised with Apptainer to guarantee a reproducible
environment across nodes; the container image bundles PyTorch~\cite{Paszke19},
PyTorch Lightning, \code{segmentation\_models\_pytorch}~\cite{Iakubovskii19} and
\code{albumentations}~\cite{Buslaev20}.

A secondary, single-GPU workstation environment (NVIDIA RTX~3050 Laptop, 4~GB) was
used for development, defect analysis and reproduction work.
Appendix~\ref{app:reproduction} documents that environment and, importantly, the
respects in which results obtained on it are \emph{not} comparable to the cluster
results reported above - most consequentially the accelerator memory limit, which
forbids the batch size of 32 used throughout the primary experiments and thereby
degrades batch-normalisation statistics.
That constraint is the reason all accuracy figures in this thesis are drawn from the
cluster.
